% -*- coding: utf-8 -*-
% !TEX root = ../main.tex

La classification de texte est traversée par un compromis entre capacité
prédictive et coût d'utilisation. À famille d'architectures et conditions
d'apprentissage comparables, augmenter le nombre de paramètres élargit en
général la capacité d'un modèle à représenter des régularités complexes. Cette
augmentation ne garantit pas à elle seule une meilleure performance : son
intérêt dépend notamment des données disponibles et de l'optimisation. Elle a
néanmoins accompagné une grande partie des progrès récents, depuis les
représentations distribuées jusqu'aux modèles contextuels préentraînés.

Cette capacité supplémentaire a un coût. Un modèle plus volumineux occupe
davantage de mémoire et demande généralement plus d'opérations pour produire
une prédiction. Or l'entraînement peut n'être réalisé qu'une fois, tandis que
l'inférence est répétée pour chaque nouveau texte. Le coût pertinent pour le
déploiement n'est donc pas seulement celui de la construction du modèle, mais
aussi celui de son stockage et de son exécution répétée. La compression cherche
à réduire ce coût récurrent tout en limitant la dégradation de la performance.
Les travaux d'élagage s'inscrivent historiquement dans cette recherche d'une
capacité plus compacte \cite{lecun1989obd,han2015learning}.

Avant d'étudier cette compression, il faut préciser ce qui est effectivement
exécuté lors d'une classification. Le traitement comporte deux problèmes
successifs : transformer une suite de symboles, de longueur variable, en une
représentation numérique exploitable, puis prendre une décision à partir de
cette représentation. Dans un système neuronal, ces fonctions sont
généralement assurées par un \emph{backbone} de représentation et par une tête
de classification. Cette décomposition fournit le fil conducteur de l'état de
l'art : comprendre l'évolution de ces deux composants, les gains obtenus et les
coûts introduits, avant de présenter les méthodes proposées pour les réduire.

\subsection{Backbone de représentation et tête de classification}
\label{ssec:etat_art_architecture}

Le backbone reçoit le texte brut et produit les caractéristiques utilisées par
le classifieur. Il comprend d'abord une étape de tokenisation, qui découpe le
texte en unités et associe un identifiant à chacune d'elles. Il contient ensuite
un mécanisme de représentation. Dans un modèle fondé sur des embeddings, chaque
identifiant est associé à un vecteur dense de petite dimension appris à partir
des données. Le terme \emph{embedding} désigne à la fois cette association et
le vecteur obtenu. Enfin, lorsque le classifieur attend un seul vecteur par
texte, le backbone doit agréger les représentations des différentes unités, par
exemple au moyen d'une moyenne ou d'un réseau qui tient compte de leur ordre.

La tête de classification reçoit la représentation produite par le backbone et
la transforme en scores de classes. Elle peut être linéaire ou composée de
plusieurs couches. Dans le cas de l'analyse de sentiments binaire, elle doit
estimer si les caractéristiques du texte sont davantage associées à un
sentiment positif ou négatif. Son apprentissage utilise donc les textes
représentés par le backbone et leurs étiquettes, alors que le backbone peut être
préentraîné sur une tâche différente et éventuellement sans annotations de
sentiment.

Cette séparation permet aussi de distinguer le coût d'apprentissage du coût
d'inférence. Un backbone peut être préentraîné une fois, puis figé lorsque le
système est déployé. Classifier un nouveau texte ne demande alors pas de
réapprendre ses paramètres : le texte est propagé dans le backbone, puis dans
la tête, avec les paramètres déjà acquis. Le préentraînement constitue donc un
coût initial, tandis que les calculs d'inférence sont répétés à chaque
prédiction.

Le backbone et la tête peuvent tous deux être comprimés, mais ces deux choix
n'ont pas la même portée. Modifier le backbone peut réduire une part importante
du coût total, au risque de changer toutes les représentations fournies au
classifieur. Agir uniquement sur la tête laisse ces représentations inchangées
et isole le coût de la décision finale. Dans les deux cas, le fait qu'un
composant soit préentraîné ne le rend pas gratuit à l'inférence : il doit encore
être stocké et exécuté pour traiter une nouvelle entrée.

\subsection{Évolution des représentations et des classifieurs de texte}
\label{ssec:etat_art_representations}

Un algorithme d'apprentissage ne traite pas directement des mots ou des
caractères : il traite des nombres. La difficulté ne consiste donc pas seulement
à attribuer un identifiant à chaque unité, mais à construire une représentation
qui conserve les informations utiles à la tâche. Les méthodes successives se
distinguent par la nature de cette information et par le coût nécessaire pour
l'obtenir.

\subsubsection{Comptages, sacs de mots et séquences locales}

Une représentation simple est le sac de mots. Chaque coordonnée du vecteur
correspond à un élément du vocabulaire et contient son nombre d'occurrences, sa
présence ou un score TF--IDF, qui pondère la fréquence locale d'une unité par sa
rareté dans le corpus. Le texte est alors décrit par les unités qu'il contient,
indépendamment de leur position. Cette représentation est
facile à calculer et à interpréter, mais elle produit des vecteurs de la taille
du vocabulaire, généralement très creux. Surtout, elle ne porte aucune notion
de proximité : \emph{excellent} et \emph{remarquable} occupent deux coordonnées
indépendantes, même si ces mots peuvent exprimer une appréciation voisine.

Le sac de mots perd également l'ordre. Pour réintroduire un contexte local, il
est possible d'utiliser des $n$-grammes. Un $n$-gramme est une suite contiguë de
$n$ unités : un unigramme contient une unité, un bigramme en contient deux et
un trigramme trois. Les bigrammes peuvent par exemple distinguer la présence de
\emph{bon} de celle de \emph{pas bon}. Leur pouvoir descriptif a toutefois un
coût : le nombre de suites possibles augmente rapidement avec $n$ et avec la
taille du vocabulaire. La majorité des combinaisons sont rares ou absentes du
corpus, ce qui limite la généralisation à des formulations nouvelles.

Les travaux de Shannon fournissent le cadre probabiliste permettant de raisonner
sur ces séquences. Ils formalisent la source de symboles, l'entropie et la
dépendance statistique \cite{shannon1948communication}. Le
traité de Gallager systématise ensuite ces outils de théorie de l'information
\cite{gallager1968information}. Ces références ne proposent pas encore de
représentation sémantique des mots. Elles établissent plutôt les concepts qui
permettent de mesurer l'incertitude d'une séquence et l'information partagée par
ses éléments.

L'analyse distributionnelle apporte un changement de perspective. Harris
observe que le comportement d'une unité linguistique peut être étudié à partir
des environnements dans lesquels elle apparaît
\cite{harris1954distributional}. Autrement dit, les contextes d'un mot donnent
une information sur son rôle dans la langue. Cette idée permet de dépasser une
représentation où chaque mot est totalement indépendant des autres, sans encore
construire de vecteurs continus.

Brown et ses coauteurs exploitent cette information de contexte dans des
modèles de langage fondés sur des classes de mots
\cite{brown1992class}. Leur méthode regroupe les mots selon leurs statistiques
de cooccurrence et utilise ces groupes dans un modèle $n$-gramme. Les classes
obtenues peuvent refléter des régularités syntaxiques ou sémantiques, et elles
réduisent la rareté par rapport à un modèle où chaque mot est traité isolément.
La similarité reste cependant discrète : deux mots appartiennent ou non à la
même classe, sans que le modèle exprime différents degrés de proximité.

À ce stade, le backbone produit donc le plus souvent un vecteur de comptages,
éventuellement enrichi par des $n$-grammes, et la tête repose sur un algorithme
d'apprentissage classique. Le classifieur naïf bayésien estime la probabilité
d'une classe en supposant les caractéristiques conditionnellement
indépendantes. Le modèle de maximum d'entropie apprend directement une
distribution sur les classes sans cette hypothèse, tandis que la machine à
vecteurs de support recherche une frontière qui sépare les exemples avec une
marge aussi large que possible. Pang, Lee et Vaithyanathan comparent ces trois
familles sur la polarité de critiques de films
\cite{pang2002sentiment}. Leur étude confirme leur intérêt pour l'analyse de
sentiments, mais montre aussi que cette tâche est plus délicate que la
catégorisation thématique : la présence d'un mot positif ou négatif ne suffit
pas toujours, car sa portée, la négation et la construction de la phrase peuvent
en inverser l'interprétation.

Ces méthodes fournissent des bases solides et interprétables, mais elles restent
fortement liées aux formes effectivement observées. Deux expressions proches
partagent peu d'information si elles n'emploient pas les mêmes unités ou les
mêmes $n$-grammes. Il devient donc nécessaire de construire une représentation
qui encode des degrés de similarité plutôt que de simples identités.

\subsubsection{Passage aux représentations distribuées}

Bengio et ses coauteurs remplacent ces classes discrètes par des
représentations distribuées continues apprises avec un modèle de langage
neuronal \cite{bengio2003neural}. Leur point de départ est la malédiction de la
dimension : le nombre de séquences possibles est immense par rapport au nombre
de séquences réellement observées. Un modèle fondé sur l'appariement exact ne
sait donc pas bien traiter une combinaison nouvelle. Dans un espace continu,
une séquence inconnue peut au contraire bénéficier de ce qui a été appris sur
des séquences composées de mots proches. Si \emph{excellent} et
\emph{remarquable} apparaissent dans des environnements comparables, leurs
vecteurs peuvent ainsi devenir voisins et partager une partie de l'information
apprise. Le modèle apprend simultanément les vecteurs des mots et la fonction
qui prédit le mot suivant.

Cette continuité améliore la généralisation, mais introduit un nouveau coût :
pour chaque contexte, le réseau doit calculer une probabilité sur l'ensemble du
vocabulaire. Morin et Bengio proposent une
décomposition hiérarchique de cette distribution
\cite{morin2005hierarchical}. Le mot cible est atteint au terme d'une suite de
décisions dans une hiérarchie, ce qui réduit fortement le nombre de sorties à
évaluer. Leur contribution porte donc moins sur la définition de l'embedding
que sur la possibilité de l'apprendre à plus grande échelle.

Le modèle de Bengio utilise encore un contexte de longueur fixée avant
l'entraînement. Mikolov et ses coauteurs montrent qu'un réseau récurrent peut
conserver un état dépendant des mots précédents et exploiter ainsi une histoire
plus longue \cite{mikolov2010rnn}. Cette approche améliore la modélisation de la
séquence, mais son entraînement demeure relativement coûteux, notamment parce
que les états successifs sont dépendants les uns des autres. Si l'objectif
principal est d'obtenir de bons vecteurs de mots, il devient alors intéressant
de simplifier le modèle de langage lui-même.

\subsubsection{Word2Vec : apprendre efficacement à partir du contexte}

Word2Vec désigne une famille de modèles légers conçus pour apprendre des
représentations distribuées à partir des contextes observés dans un grand
corpus. Mikolov et ses coauteurs introduisent deux architectures
\cite{mikolov2013efficient}. Le modèle \emph{Continuous Bag-of-Words} (CBOW)
utilise les unités voisines pour prédire l'unité centrale. Skip-gram effectue la
tâche inverse : il part de l'unité centrale et cherche à prédire son contexte.
Par exemple, CBOW peut utiliser le contexte \emph{ce film est vraiment …} pour
prédire une unité telle que \emph{excellent}, alors que Skip-gram part de
\emph{excellent} pour apprendre à reconnaître les unités susceptibles de
l'entourer. Dans les deux cas, la tâche de prédiction sert de signal
d'apprentissage ; les vecteurs contenus dans la couche d'embedding sont ensuite
réutilisés dans d'autres applications.

Cette simplification ne résout pas entièrement le coût associé à un grand
vocabulaire. Mikolov et ses coauteurs proposent alors l'échantillonnage négatif,
ou \emph{negative sampling} \cite{mikolov2013distributed}. Plutôt que de
comparer systématiquement le mot cible à tous les éléments du vocabulaire,
l'entraînement oppose une association réellement observée à un petit nombre
d'associations artificielles. Le calcul devient beaucoup moins coûteux. Le
sous-échantillonnage des mots très fréquents évite en outre qu'ils dominent
l'apprentissage alors qu'ils sont souvent peu discriminants. Goldberg et Levy
reprennent ensuite cet objectif et en donnent une dérivation plus explicite
pour Skip-gram \cite{goldberg2014word2vec}. Leur analyse aide à comprendre que
l'échantillonnage négatif est un objectif d'apprentissage de représentations à
part entière, et pas uniquement une implémentation accélérée d'un modèle de
langage complet.

\subsubsection{Embeddings statiques et prise en compte de l'ordre}

Un embedding ainsi appris est dense et rapproche les unités apparaissant dans
des contextes similaires. Pour classifier un texte, la solution la plus simple
consiste à agréger ses embeddings, puis à fournir le vecteur obtenu à une tête
neuronale, par exemple un perceptron multicouche. Cette architecture est rapide
et sépare clairement le backbone de la tête. Elle ne peut toutefois pas
récupérer l'ordre perdu pendant l'agrégation : approfondir uniquement la tête ne
corrige donc pas les limites de la représentation qui lui est fournie.

Les architectures suivantes déplacent la modélisation de la séquence dans le
backbone. Kim applique des convolutions à la suite des vecteurs de mots afin de
détecter des motifs locaux utiles à la classification
\cite{kim2014cnn}. Les réseaux récurrents traitent plutôt les unités
successivement et conservent un état qui résume le contexte précédent. Ces deux
familles exploitent mieux l'ordre que la moyenne d'embeddings : les convolutions
privilégient les motifs locaux, tandis que les récurrences peuvent en principe
relier des positions plus éloignées. Elles couplent cependant davantage la
représentation et la classification, et leur entraînement est plus coûteux que
celui d'une simple tête appliquée à un vecteur fixe.

\subsubsection{Du mot aux unités sous-lexicales}

Les modèles qui associent un vecteur à chaque mot complet rencontrent une autre
limite : leur vocabulaire est nécessairement fini. Une forme rare ou absente du
corpus d'apprentissage ne possède donc pas de représentation fiable. Les unités
sous-lexicales réduisent ce problème en représentant un mot à partir de fragments
plus fréquents. Des formes comme \emph{jouer}, \emph{joueur} et \emph{rejouer}
peuvent alors partager certains fragments, même si la segmentation apprise ne
coïncide pas toujours avec un découpage morphologique.

Le \gls{bpe} construit ce vocabulaire en fusionnant progressivement les paires
de symboles les plus fréquentes. Sennrich, Haddow et Birch montrent l'intérêt de
ce principe pour le traitement des mots rares en traduction neuronale
\cite{sennrich2016subword}. Le recours à des sous-unités limite le nombre de
mots entièrement inconnus, mais allonge les séquences et ne résout pas à lui
seul la dépendance au contexte.

\subsubsection{Vers des représentations contextuelles}

Les embeddings statiques conservent en outre une limite fondamentale : une même
unité reçoit le même vecteur dans tous ses contextes et représente donc mal la
polysémie. Le mot \emph{avocat}, par exemple, conserve une représentation unique
dans \emph{l'avocat prépare sa plaidoirie} et \emph{elle mange un avocat}, alors
que son sens change complètement. Les transformeurs préentraînés répondent à
cette limite en produisant des représentations contextuelles. BERT construit
ainsi le vecteur d'une unité à partir de son contexte gauche et droit, puis
l'ensemble du modèle peut être ajusté pour une tâche de classification avec une
couche de sortie supplémentaire \cite{devlin2019bert}. Cette capacité améliore
la prise en compte du contexte, mais elle repose sur un backbone beaucoup plus
volumineux et plus coûteux à l'inférence.

\subsection{Pourquoi comprimer le modèle obtenu ?}
\label{ssec:etat_art_compression}

Le besoin de compression posé en ouverture ne vient pas d'un réentraînement du
modèle à chaque requête : une fois déployés, le backbone et la tête conservent
leurs paramètres appris. Il vient du coût répété de leur utilisation. Chaque
nouveau texte doit être représenté, puis propagé dans le classifieur. Selon le
périmètre considéré, ce coût comprend la taille du modèle stocké, la mémoire
mobilisée, le nombre d'opérations, les accès mémoire, la latence et la
consommation énergétique.

Une méthode de compression doit préciser le composant et la ressource qu'elle
cherche à réduire. Une diminution théorique du nombre de paramètres ne garantit
pas, à elle seule, une diminution de la latence sur le matériel utilisé. Les
formats de stockage, les accès mémoire et la capacité des bibliothèques à
exploiter la structure obtenue déterminent également le gain réellement
observable.

Plusieurs familles de méthodes répondent à ce problème. La quantification réduit
la précision numérique des poids ou des activations. Les factorisations de
faible rang remplacent de grandes matrices par des produits de matrices plus
petites. La distillation entraîne une nouvelle architecture compacte à imiter
un modèle plus puissant. Le pruning conserve l'architecture de départ mais en
retire les paramètres ou les structures jugés les moins utiles. Ces approches
peuvent être combinées : Tu et ses coauteurs associent par exemple un critère de
pruning fondé sur l'information de Fisher à une quantification non uniforme
\cite{tu2016fisher}.

\subsubsection{Distillation de connaissances}

La distillation formalise une relation entre un modèle enseignant et un modèle
étudiant. Hinton, Vinyals et Dean proposent de ne pas entraîner l'étudiant
uniquement avec les étiquettes finales, mais aussi avec les probabilités
produites par l'enseignant \cite{hinton2015distilling}. Ces cibles souples
indiquent quelles classes l'enseignant considère comme proches ou facilement
confondables. Elles transmettent ainsi une information absente d'une étiquette
dure. Une température contrôle leur lissage, et l'entraînement combine
généralement l'imitation de l'enseignant avec la prédiction des vraies
étiquettes.

La limite de cette formulation initiale est qu'elle exploite surtout la sortie
finale. Les travaux ultérieurs transfèrent également des représentations
intermédiaires ou des relations entre exemples. La synthèse de Gou et ses
coauteurs organise ces méthodes selon le type de connaissance transférée et
selon le scénario d'entraînement : enseignant préentraîné et figé, apprentissage
conjoint ou auto-distillation \cite{gou2021survey}. Une synthèse plus récente
étend cette lecture aux transformeurs, aux modèles de fondation et aux grands
modèles de langage, tout en soulignant les difficultés créées par l'écart de
capacité et de structure entre enseignant et étudiant
\cite{mansourian2025survey}.

\subsubsection{Pruning structuré et non structuré}

Le pruning non structuré retire des connexions individuelles. Sa granularité
fine permet d'atteindre une forte sparsité, c'est-à-dire une forte proportion de
poids annulés, mais les zéros restent stockés et calculés si les matrices
conservent un format dense. Un gain physique exige
alors un format creux et des noyaux capables de l'exploiter. Le pruning
structuré retire des unités complètes, comme des neurones, des canaux ou des
blocs. Il réduit plus facilement les dimensions réelles des matrices, mais son
choix est plus contraint et peut dégrader plus rapidement la performance.

Dans les deux cas, le problème essentiel est le même : il faut estimer
l'importance d'un élément avant de le supprimer. Le critère exact consisterait
à retirer chaque candidat, à réévaluer tout le modèle puis à le restaurer. Ce
calcul est trop coûteux ; les travaux de la littérature proposent donc des
approximations fondées successivement sur la courbure, la magnitude, les
gradients ou les activations.

\subsection{Évolution des critères de pruning}
\label{ssec:etat_art_pruning}

\subsubsection{De la courbure aux méthodes pratiques}

Optimal Brain Damage, proposé par LeCun, Denker et Solla, fait partie des
premières méthodes qui donnent une définition analytique de l'importance d'un
poids \cite{lecun1989obd}. L'idée consiste à approximer l'augmentation de la
fonction de coût provoquée par sa suppression. Le score dépend à la fois de la
valeur du poids et de la courbure locale de la fonction de coût. Pour rendre le
calcul possible, la méthode ne conserve que la diagonale de la Hessienne et
suppose que le réseau se trouve près d'un minimum. Elle apporte donc plus
d'information qu'un classement fondé sur la seule magnitude, mais néglige les
interactions entre paramètres.

Optimal Brain Surgeon est introduit par Hassibi et Stork pour lever précisément
cette hypothèse \cite{hassibi1993obs}. La méthode utilise l'inverse de la
Hessienne complète et ajuste les poids conservés afin de compenser la connexion
supprimée. Elle modélise mieux les dépendances entre paramètres et peut éviter
certaines suppressions incorrectes d'OBD. Son coût mémoire et calculatoire
augmente toutefois très rapidement avec le nombre de poids. Cette limite rend
la méthode difficile à appliquer telle quelle aux réseaux modernes et explique
l'intérêt persistant pour des approximations diagonales.

Avec l'augmentation de la taille des réseaux, Han et ses coauteurs remettent en
avant un critère beaucoup plus simple : la magnitude des poids
\cite{han2015learning}. Leur contribution principale ne réside pas uniquement
dans le score, mais dans un protocole en trois étapes : entraînement dense,
suppression des petites connexions, puis réentraînement des poids restants. Ils
montrent que cette dernière étape est déterminante et que plusieurs cycles
progressifs sont préférables à une suppression brutale. La méthode passe à
l'échelle, mais la magnitude ne mesure pas directement l'effet d'un poids sur
la prédiction et la sparsité irrégulière n'accélère pas automatiquement le
matériel généraliste.

\subsubsection{Information de Fisher et critères du premier ordre}

Tu et ses coauteurs proposent d'utiliser l'information de Fisher pour mesurer
la sensibilité de la distribution de sortie aux paramètres
\cite{tu2016fisher}. La diagonale de cette matrice peut être estimée à partir des
gradients accumulés pendant l'entraînement. Les paramètres associés à une faible
information sont considérés comme moins importants. Les auteurs constatent
cependant que les très petites valeurs sont difficiles à estimer de manière
fiable et combinent donc l'information de Fisher avec la magnitude. Leur étude
montre aussi que pruning et quantification peuvent être pensés conjointement.
La validation initiale porte néanmoins sur un réseau convolutif et le jeu MNIST,
ce qui impose de vérifier la transposition au texte.

Molchanov et ses coauteurs cherchent ensuite à approximer plus directement la
variation de la fonction de coût à l'aide d'un développement de Taylor du
premier ordre \cite{molchanov2017pruning}. Le score combine l'activation d'un
élément et le gradient qui lui est associé. Il évite le calcul d'une Hessienne
et peut être appliqué à des structures telles que les filtres convolutifs. Dans
leurs expériences, cette approximation offre un compromis favorable entre
qualité du classement et coût de calcul. Elle suppose toutefois l'accès à des
données et, lorsque le gradient provient de la fonction supervisée, à leurs
étiquettes.

Cette dépendance aux labels motive les travaux de Sapkota et Bhattarai
\cite{sapkota2023random}. Les auteurs remplacent le gradient supervisé de
sortie par un gradient aléatoire normalisé, puis propagent ce signal dans le
réseau pour estimer l'importance des neurones. Le pruning peut ainsi utiliser
des données sans annotations. Les résultats sont toutefois obtenus sur des
réseaux de vision et publiés dans un article de workshop ; leur transposition à
la classification de texte reste donc à évaluer.

\subsubsection{Dépendance entre couches et information mutuelle}

Les méthodes précédentes évaluent surtout la sensibilité de la fonction de coût
à un poids ou à une activation. MINT, proposé par Ganesh, Corso et Yasaei Sekeh,
change de point de vue en mesurant l'information transmise entre deux couches
successives \cite{ganesh2020mint}. Un groupe d'unités est considéré comme
important s'il apporte une information qui n'est pas déjà contenue dans les
autres activations de la couche. La méthode vise ainsi à conserver les
connexions qui participent le plus à la dépendance entre les représentations
amont et aval.

L'estimation directe de l'information mutuelle est difficile dans un espace de
grande dimension. MINT s'appuie donc sur un estimateur géométrique qui évite de
reconstruire explicitement les densités. Les travaux de Yasaei Sekeh et Hero
fondent cette estimation sur une comparaison géométrique de populations au
moyen d'un arbre couvrant minimal \cite{yasaei2019gmi}. La construction d'une
population représentant l'indépendance conditionnelle utilise un bootstrap par
plus proches voisins, relié aux travaux de Sen et ses coauteurs sur les tests
d'indépendance conditionnelle \cite{sen2017ci}.

Cette approche tient compte des activations et des relations entre couches,
mais elle introduit de nouveaux choix : nombre d'exemples de calibration,
constitution des groupes, seuil de suppression et protection des couches. Le
calcul de nombreux scores géométriques peut également être coûteux. L'article
original évalue principalement des architectures de vision, même s'il inclut
un perceptron multicouche sur MNIST. La capacité de la méthode à se transposer à
une tête de classification de texte n'est donc pas directement établie par les
expériences publiées.

Les travaux récents prolongent cette idée selon des objectifs différents. MIPP
cherche à préserver l'information mutuelle entre les activations de couches
adjacentes et relie cette conservation à la capacité de réentraîner le réseau
après un pruning structuré \cite{westphal2025mipp}. Huang et Pao étudient plutôt
la redondance interne des couches de représentation, mesurée par
l'information mutuelle et la corrélation totale
\cite{huang2025interpretable}. Leur méthode peut décider du pruning sans labels
une fois le modèle entraîné. Ces contributions montrent que l'expression
« pruning par information mutuelle » recouvre plusieurs buts : conserver
l'information transmise, éliminer la redondance ou préserver l'information
utile à la sortie. Ces objectifs ne conduisent pas nécessairement au même
masque.

\subsection{Positionnement et périmètre du travail}
\label{ssec:etat_art_positionnement}

Au terme de cette revue, deux constats orientent le travail. D'une part, une
représentation contextuelle riche peut améliorer la classification, mais elle
rend plus difficile l'étude isolée du classifieur. D'autre part, les critères de
pruning forment une progression allant de scores locaux et peu coûteux vers des
scores qui utilisent davantage d'information sur le comportement du réseau. La
magnitude est simple mais indirecte ; OBD ajoute la courbure de la fonction de
coût, tandis que MINT évalue la dépendance entre représentations successives.

Le pipeline retenu privilégie donc une séparation explicite entre
représentation et décision. Le texte est tokenisé en unités sous-lexicales avec
le \gls{bpe}. Un modèle CBOW apprend un embedding de dimension 100 pour chaque
unité, puis le backbone représente le texte par la moyenne de ces embeddings.
Le vecteur obtenu est enfin transmis à une tête de classification multicouche.
Cette moyenne fournit une représentation compacte et de taille fixe, mais elle
perd l'ordre des unités : elle peut par exemple rapprocher \emph{ce film est
bon} et \emph{ce film n'est pas bon} alors que leur polarité diffère.

Le modèle d'embedding est préentraîné une fois, puis figé. Il n'est donc pas
réentraîné lors de chaque classification. Il reste néanmoins exécuté pour les
nouveaux textes, à travers la tokenisation, la consultation des embeddings et
leur agrégation. Le backbone est conservé identique et n'est pas pruné dans les
expériences : la variable étudiée est uniquement la tête. Ce choix maintient les
représentations d'entrée constantes et permet d'attribuer les différences
observées au critère de compression. Il impose en contrepartie de distinguer les
gains mesurés sur la tête de ceux du pipeline complet.

Le protocole compare principalement OBD et MINT sur la même tête dense initiale.
Ces méthodes représentent deux conceptions différentes de l'importance : OBD
estime l'effet local de la suppression d'un poids à partir de la courbure, alors
que MINT s'appuie sur l'information transmise entre des activations. Le modèle
est évalué immédiatement après l'élagage, puis après un réentraînement contraint
par le masque. Une petite architecture dense entraînée directement fournit un
point de comparaison à nombre de poids voisin.

La comparaison porte sur des modèles produits à un niveau de \gls{sparsity}
comparable, et non sur les valeurs brutes des scores, qui n'ont ni la même
échelle ni la même granularité. L'application de MINT à une tête de
classification de sentiments constitue une adaptation de travaux surtout
évalués en vision. La distillation et le pruning sans annotations représentent
des prolongements pertinents, mais ne font pas partie de la comparaison
principale. Les conclusions concernent donc l'élagage d'une tête multicouche
alimentée par des embeddings statiques ; elles ne prétendent pas désigner une
méthode universellement optimale pour tous les modèles de \gls{nlp}.
